% Source template: analyze_iq_steps.py copies this file into the output directory.
% Compile the generated copy there: pdflatex -interaction=nonstopmode -halt-on-error step_report.tex
\documentclass[a4paper,11pt]{article}
\usepackage[margin=20mm]{geometry}
\usepackage[T1]{fontenc}
\usepackage{lmodern}
\usepackage{amsmath,booktabs,pgfplots,pgfplotstable,siunitx}
\usepackage[hidelinks]{hyperref}
\pgfplotsset{compat=1.18}
\definecolor{uzblue}{HTML}{1F77B4}
\definecolor{uzorange}{HTML}{FF7F0E}
\definecolor{uzgreen}{HTML}{2CA02C}
\definecolor{uzred}{HTML}{D62728}
\definecolor{uzpurple}{HTML}{9467BD}
\setlength{\parindent}{0pt}
\setlength{\parskip}{5pt}
\hypersetup{pdftitle={Experimental Evaluation of PMSM q-Axis Current Control}}
\begin{document}
\begin{center}
{\Large\bfseries Experimental Evaluation of PMSM\\q-Axis Current Control}\par
\smallskip
{\large Influence of Proportional Gain on Transient Response}
\end{center}

\section*{1. Objective and experimental setup}
A PMSM in controller-in-the-loop (CIL) operation is controlled by a parallel PI
q-axis current controller. Only $K_p$ is varied; $K_i$ remains constant. The nominal
design point is $K_p=8.8\,\mathrm{V/A}$ with an approximate design bandwidth of
$800\,\mathrm{Hz}$; this is not an experimentally measured bandwidth. The numerical
value of $K_i$ is not recorded. The experiment is executed automatically on
UltraZohm and logged using JavaScope. CIL operation and the nominal design are
specified by the experiment description.

The recorded reference changes from 0 to 1\,A (amplitude 1\,A). The 5 evaluated gains are 2, 5, 8.8, 12, 15\,V/A. The recorded motor speed was approximately 200\,rpm. Timestamp spacing is 0.1\,ms, corresponding to a logging rate of 10000\,Hz.

\section*{2. Experimental results}
\begin{center}
\begin{tikzpicture}
\begin{axis}[
width=0.97\linewidth,height=9.5cm,
xmin=-2, xmax=15,
xlabel={Time relative to reference step / ms},
ylabel={q-axis current / A},
grid=major,grid style={gray!20},
unbounded coords=jump,enlarge y limits=0.12,
tick label style={font=\small},
legend style={font=\small,draw=none,at={(0.5,-0.17)},anchor=north},
legend columns=3,
]
\addplot[very thick,solid,uzblue] table[x=time_ms,y=iq_kp_2,col sep=comma]{step_plot_data.csv};
\addlegendentry{$K_p=2$\,V/A}
\addplot[very thick,solid,uzorange] table[x=time_ms,y=iq_kp_5,col sep=comma]{step_plot_data.csv};
\addlegendentry{$K_p=5$\,V/A}
\addplot[very thick,solid,uzgreen] table[x=time_ms,y=iq_kp_8p8,col sep=comma]{step_plot_data.csv};
\addlegendentry{$K_p=8.8$\,V/A}
\addplot[very thick,solid,uzred] table[x=time_ms,y=iq_kp_12,col sep=comma]{step_plot_data.csv};
\addlegendentry{$K_p=12$\,V/A}
\addplot[very thick,solid,uzpurple] table[x=time_ms,y=iq_kp_15,col sep=comma]{step_plot_data.csv};
\addlegendentry{$K_p=15$\,V/A}
\addplot[black!65,dashed,thin,const plot] table[x=time_ms,y=iq_ref,col sep=comma]{step_plot_data.csv};
\addlegendentry{Reference}
\end{axis}
\end{tikzpicture}
\end{center}
\vspace{3mm}
{\small\textbf{Figure 1.} Measured q-axis current responses for different proportional
gains. All curves are aligned to the reference step at $t=0$.}

\newpage
\section*{3. Quantitative evaluation}
\begin{center}
\pgfplotstabletypeset[
col sep=comma,
columns={kp_v_per_a,rise_time_ms,overshoot_percent,settling_time_ms,iae_a_ms},
columns/kp_v_per_a/.style={column name={$K_p$ / (V/A)},fixed,precision=1},
columns/rise_time_ms/.style={column name={$t_r$ / ms},fixed,fixed zerofill,precision=3},
columns/overshoot_percent/.style={column name={Overshoot / \%},fixed,fixed zerofill,precision=2},
columns/settling_time_ms/.style={column name={$t_s$ / ms},fixed,fixed zerofill,precision=2},
columns/iae_a_ms/.style={column name={IAE / ($\mathrm{A}\cdot\mathrm{ms}$)},fixed,fixed zerofill,precision=2},
every head row/.style={before row=\toprule,after row=\midrule},
every last row/.style={after row=\bottomrule},
]{step_metrics.csv}
\end{center}

Rise time $t_r$ uses the first 10--90\% crossings of the commanded step amplitude,
with linear interpolation between raw samples. Settling time $t_s$ is the first
sample after the last excursion outside the $\pm2\%$ band around the final
reference, maintained through the remainder of HOLD. IAE is the trapezoidal
integral of $|i_q^\star-i_q|$ from the reference edge to the last HOLD sample.
Metrics use unfiltered data and recorded timestamps; NaN denotes an unavailable
metric rather than zero.

\smallskip
\textbf{Voltage utilization.}
Maximum recorded $|u_q^\star|$ during HOLD is 18.77\,V. The maximum command-vector magnitude $\sqrt{(u_d^\star)^2+(u_q^\star)^2}$ is 18.77\,V. Using externally supplied $U_{\mathrm{dc}}=48$\,V (not recorded), the nominal linear-SVPWM bound is 27.71\,V. Maximum vector utilization is 67.7\,\%. The recorded command magnitude remains below this bound. This compares voltage commands, not measured terminal voltage.

\section*{4. Discussion}
The smallest overshoot is 7.29\,\% at $K_p=5$\,V/A. The shortest settling time is 1.10\,ms at $K_p=8.8$\,V/A. The lowest IAE is 0.29\,A\,ms at $K_p=8.8$\,V/A. The slowest rise occurs at $K_p=2$\,V/A (0.741\,ms). Every tested gain above the nominal value reduces rise time but increases overshoot. Alternating post-rise excursions beyond the $\pm2\%$ band indicate ringing; the largest count is 13 side changes at $K_p=15$\,V/A. These trade-offs apply to the tested operating point, not a universal optimum. Neither rise time nor the nominal 800 Hz design target establishes a measured closed-loop bandwidth. The waveform alone does not identify the cause of ringing.

\section*{5. Reproducibility}
\textbf{Measurement type:} CIL (experiment description).\\
\textbf{Original JavaScope recording:}
\texttt{Log\_2026-11-04\_20-03-07.csv}

Operating conditions and logging rate are summarized in Section~1; the DC-link
source is identified with the voltage-utilization assessment. The table reads
\texttt{step\_metrics.csv}; PGFPlots reads \texttt{step\_plot\_data.csv}.
Only the figure is resampled by linear interpolation within contiguous recorded
intervals; missing coverage is omitted. The reference retains its discrete step.
Numerical metrics use the original complete HOLD intervals. Incomplete or invalid
measurements are excluded with analysis warnings; logged DUT violations are
flagged in the discussion.
\end{document}
